\subsection{Learned-model evaluation protocol}\label{app:pairs-v102}
Six objectives continue each of five common conservative models for
1,200 Adam updates, with learning rate $3\times10^{-5}$, flow batch 32
and query batch eight. The flow data and seed-specific minibatch
sequences are shared. Each of six queries has 64 training and 16
validation endpoint pairs. Centers are sampled coordinatewise from
$[0.65,1.3]$ times baseline demand (seed 102111); half-widths 2.5\% and
15\% occur equally often. Only one demand changes per pair, while
background routes remain responsive.

The flow loss averages $[(f_e^L-F_e^S)/s_e]^2$ over minibatch states
and edges, with $s_e=0.01\sqrt{\operatorname{mean}_{\rm train}[(F_e^S)^2]}$ from the
2,500 training states. Each auxiliary objective adds one term
$\lambda\,\mathcal L_{\mathrm{aux}}$ to this flow loss, with the weight
grid $\lambda\in\{0.001,0.003,0.01,0.03,0.1,0.3,1\}$ searched
identically for every objective. Each normalizer carries the units of
its own target, so that the common grid explores comparable
regularization ranges: rescaling the outcome $T\mapsto cT$ rescales
the weighted terms' normalizers with them and leaves every effective
weight unchanged. With
$\langle\cdot\rangle_j$ denoting the average over query $j$'s training
centers or endpoint labels, the training-only scales are
\[
\begin{aligned}
u_j&=\max\{\sqrt{\langle\|J_j^S\|_2^2\rangle_j},10^{-6}\},&
v_j&=\max\{\sqrt{\langle\|D^SJ_j^S\|_2^2\rangle_j},10^{-6}\},&
b_j&=\max\{\sqrt{\langle\|z_j^S\|_2^2\rangle_j},10^{-6}\},\\
a_j&=\max\{\sqrt{\langle\mathcal K_j(z_j^S)^2\rangle_j},10^{-6}\},&
\bar a_j&=\max\{\sqrt{\langle\mathcal K_j(D^SJ_j^S)^2\rangle_j},10^{-6}\}.
\end{aligned}
\]
The auxiliary terms are: unweighted response
$\langle\|J_j^L-J_j^S\|_2^2/u_j^2\rangle$; weighted response
$\langle\|D^LJ_j^L-D^SJ_j^S\|_2^2/v_j^2\rangle$, with
$D^r=\operatorname{diag}(\nabla T(F^r))$ the outcome-gradient weights
of Eq.~\eqref{eq:Cdef} evaluated at each model's own flows (TSTT:
$\nabla_eT=t_e+F_et'_e(F_e)=t_{0,e}[1+0.75(F_e/c_e)^4]$);
the pointwise boundary seminorm
$\langle[\mathcal K_j(D^LJ_j^L-D^SJ_j^S)/\bar a_j]^2\rangle$; the
endpoint-edge term
$\langle\|z_j^L-z_j^S\|_2^2/b_j^2\rangle$; and the endpoint-cut bundle,
which adds $0.01$ times
$\langle[\mathcal K_j(z_j^L-z_j^S)/a_j]^2\rangle$ to the endpoint-edge
term (a fixed $0.01$ inner ratio, not an independently tuned
objective). The three
response-family objectives share one set of derivative labels; the two
endpoint objectives share one set of endpoint labels. The weighted
objective is run at every grid weight under this own-scale
normalization from the shared initializations.

Checkpoint selection uses one uniform rule for every objective and
weight: among the common initial model (update zero) and every 100th
continuation update whose validation flow error is at most $1.1$ times
the initial model's value, take the minimum validation boundary
mean-square error, using the same 96 validation pairs throughout. The
initial model ensures the pool is never empty. The weight $\lambda^*$ of
each objective is then chosen on validation data only: among the grid
values whose seed-median selected validation flow error stays within
$1.1$ times the seed-median initial error, minimize the seed-median
validation boundary error. The flow-only control reuses the flow
continuation under the same uniform rule. Test paths never enter
selection or weight choice. The extended-budget replication of
Section~\ref{sec:pairresults} repeats the response and flow
continuations with 3{,}600 updates under this unchanged protocol
(per-network costs in Appendix~\ref{app:anaheimlearning}).

The final evaluation uses 96 fixed test paths, all generated by one
rule: for each active-query group $k\in\{1,2,3,6\}$, draw a base state
uniformly from $[0.65,1.3]$ times baseline demand, choose the $k$
active queries in round-robin order, add independent changes of
$15$--$30\%$ of their baseline demands with alternating same-direction
and mixed-direction signs, and solve all 41 states of the resulting
straight demand path. Twenty-four paths (six per group, seed 102112)
and seventy-two paths (eighteen per group, seed 102113) were generated;
models and validation choices were fixed before evaluation on the
additional paths. For every path, flow error first averages relative
Euclidean errors across its states. Attribution statistics average active
queries and seeds within each path and then take medians across the 96 paths;
seeds do not add independent paths. For column error $e_j$ and reference column
$c_j$, relative column, system and unrestricted-boundary errors divide
$\|e_j\|_2$, $|\mathbf1^\top e_j|$ and $\mathcal K_j(e_j)$ by the
corresponding reference values, with absolute floor $10^{-6}$. Geometric error
uses $\max_{U\in\mathcal F_j^{\rm geo}}|\mathbf1_U^\top Be_j|$ and
the same maximum for $c_j$. Percentages multiply ratios by 100. The
incremental-flow error at the path endpoints is
$\norm{[f^L(q^1)-f^L(q^0)]-[f^S(q^1)-f^S(q^0)]}_2/\norm{f^S(q^1)-f^S(q^0)}_2$,
reported with the absolute numerator because the reference increment
is small relative to total flows (per path--seed medians of both, and
the smallest reference increment, in
Appendix~\ref{supp:learningdetail}).

The decomposition in Eq.~\eqref{eq:coldecomp} is evaluated at the same 41 states
with the same quadrature. Component norms describe error mechanisms; cancellation
prevents treating their norms as additive shares. The fixed-screenline sign
comparison in Appendix~\ref{supp:screenlinelearned} uses the separate
pre-specified six-query baseline segment and is not pooled with these 24 paths.
Table~\ref{tab:learning-costs} reports the incremental computation
costs of training-label generation, common validation and explanation
timing under the shared CPU protocol.
\begin{table}[t]\centering\footnotesize
\caption{Incremental computation costs (seconds) under the shared
continuation protocol. Training-label and common-validation costs are
measured separately on the same CPU protocol; training times are
continuation-time medians over the weight grid and seeds, and
explanation times are sequential medians for the selected models over
the 24-path timing subset defined below (totals for that subset).}
\label{tab:learning-costs}
\begin{tabular}{lrrrr}\toprule Labels & Training labels & Common validation & Continuation & Explanation\\\midrule
Flow only & 0.00 & 6.45 & 246.0 & 7.34\\
Response family & 16.12 & 6.45 & 155--425 & 7.13\\
Endpoint (edge, cut) & 25.75 & 6.45 & 441--497 & 7.27--7.59\\
\bottomrule\end{tabular}\end{table}
Training labels comprise 768 endpoint states for either endpoint
objective, or 384 center states and their sensitivities for the three
response-family objectives (unweighted, weighted, boundary-seminorm),
which share one label set. Every method additionally uses the same 192
endpoint states for validation; unused center-validation derivatives
are excluded. Existing flow data, pretrained models and the common test
reference are excluded from these incremental costs. The 984 test
states and sensitivities took 74.4\,s in the reference run.
Explanation timing covers a 24-path timing subset of the 96 test paths
(the seed-102112 paths, six per active-query group; forward
evaluation, six directional derivatives per state, quadrature and node
aggregation, excluding model loading): the table entries are totals
over that subset, about $0.3$\,s per path, not the full 96-path
evaluation. Continuation times include
initialization, validation and saving under concurrent jobs, pooled
across the original run and the renormalized-weight re-run under
lighter load; they describe the actual training expense and do not
establish a method speed advantage.

